\documentclass[10pt,twocolumn]{article}
\usepackage[utf8]{inputenc}
\usepackage[english,russian]{babel}
\usepackage{amsmath,amssymb,mathtools,bm}
\usepackage{geometry}
\geometry{a4paper,margin=1cm}

\DeclareMathOperator{\Rank}{Rank}
\DeclareMathOperator{\Colim}{Colim}
\DeclareMathOperator{\Period}{Period}

\begin{document}

\title{\tiny\bf Унивалентные Основания Мультимасштабного Метатопоса $7n/M$}
\author{\tiny Исследовательская группа $7n$}
\date{\vspace{-9ex}}
\maketitle

\begin{abstract}
\tiny
В работе излагаются предельно обобщенные унивалентные гомотопические основания метатопоса $7n/M$. Логика, число и геометрия выводятся как фазовые состояния калибровочной системы на пересечении свойств контекста счёта и относительных единиц (о.е.).
\end{abstract}

\paragraph{Онтологический инвариант.} 
\textit{Любой элемент есть относительная единица внешней системы либо сложная внутренняя система отношений. Счёт единиц производится через отношения, а отношений — через единицы. Пространство материализуется через единицы, а Время разворачивается через отношения.}

\section*{1. Аксиома Унивалентности и Пространство $S^2$}
В основаниях метатопоса традиционные множества заменяются типами высших гомотопий. Согласно аксиоме унивалентности Воеводского, эквивалентность структур тождественна равенству:
\begin{equation}
(A \simeq B) \simeq (A = B)
\end{equation}

Непрерывный перевод логических истинностных значений в геометрические координаты на сфере $S^2$ задается связностью Мальцева $\nabla_{7n}(s, \theta)$, где широта $\theta \in [0, \pi]$ определяет масштаб и баланс термодинамических оппозиций (Энтропия $\rightleftharpoons$ Негэнтропия):
\begin{equation}
\nabla_{7n}(s, \theta) = \cos(\theta) \cdot e_0 \otimes \mathbf{\hat{\mathbf{I}}} + \frac{\sin(\theta)}{\sqrt{3}} \sum_{i,j,k=1}^{7} f_{ijk} \, e_k \otimes \left( \mathbf{\hat{\mathbf{P}}}_{i \cdot} \otimes \mathbf{\hat{\mathbf{P}}}_{\cdot j} \right)
\end{equation}

\begin{itemize}
    \item \textbf{Булев предел ($\theta = 0, \pi$):} Динамический неассоциативный член обнуляется ($\sin\theta = 0$). Пространство абсолютно, время локально заморожено. Система редуцируется в бинарный базис День/Ночь, соответствующий сужению ассоциированной матрицы плотности до тривиального ядра.
    \item \textbf{Тернарный экватор ($\theta = \pi/2$):} Статический каркас исчезает ($\cos\theta = 0$). Время абсолютно, пространственные границы стерты. Включается неассоциативное перемешивание через тензор плоскости Фано $f_{ijk}$ на лупах Муфанг.
\end{itemize}

\section*{2. Системный Мультивывод 0, 1 и Контекста}
Первичный Нуль ($\mathbf{0}$) и Единица ($\mathbf{1}$) выводятся как фазовые режимы редукции калибровочного автомата по четности шага $q$:
\begin{equation}
\mathcal{W}_{\text{inv}}^{(157)} \cdot X^{(157)} = \left[ \delta_{(q \bmod 2), 0} \mathcal{M}_{\text{cycl}}(10) + \delta_{(q \bmod 2), 1} \mathbf{0} \right] X^{(157)} \pmod{13}
\end{equation}

При четном $q$ на решетке $N(12)=157$ запускается эволюция дискретного разностного оператора $\Delta_{\tau}$:
\begin{equation}
\Delta_{\tau} X^{(157)}_{k} \equiv 10^k \cdot (!13)^{\lfloor k/6 \rfloor} X^{(157)}_* \pmod{13}
\end{equation}

Поскольку модулярный дефект субфакториала удовлетворяет сравнению $!13 \equiv -1 \pmod{13}$ (согласно аналогу теоремы Вильсона для деранжирований), а число 10 имеет мультипликативный порядок $\text{ord}_{13}(10)=6$, локальный период автомата удваивается до $\Period_{\tau} = 12$. На 12-м шаге топологическая эверсия знака полностью компенсируется:
\begin{equation}
\mathbf{\hat{T}}_{\text{inv}}^{\times 12} \equiv \mathbf{Id}_{\mathcal{C}} \implies \Rank_{\mathbb{F}_2}(\mathbf{\hat{\Xi}}^{(157)}) \equiv 1
\end{equation}
Это возвращает Объектную Единицу ($\mathbf{1}$), тогда как нечетный шаг через факторизацию в терминальный объект категории $\mathbf{1}_{\mathcal{C}}$ выводит Объектный Нуль ($\mathbf{0}$).

\section*{3. Теорема об Изоморфизме и Тавтологии}
\textbf{Теорема.} \textit{В высшей металогике $7n/M$ логическая тавтология тождественна геометрическому изоморфизму структур.}
\\
\textbf{Доказательство.} Истинность суждения не является плоскостным тождеством знака $A \equiv A$. Успешное прохождение через 12-шаговую монодромию Коловорота (включая инверсию ориентации пути на шаге $k=6$) возвращает тождественный функтор $\mathbf{Id}_{\mathcal{C}}$ над полем $\mathbb{F}_2$. Совпадение начального и конечного гомотопических типов доказывает их изоморфизм, что эквивалентно высшему выводу Истины.

\section*{4. Регистровые Резонаторы 6, 7, 10}
Масштабные фильтры счета кратны фундаментальным константам. Вероятности комбинаций упорядоченных перестановок и деранжирований $(3!) \times (4!) = 144$ за счет инцидентности Фано с тильдой $\tilde{\mathcal{F}}_{\text{ano}}$ осуществляют сплит масштаба $\pm(3-4-3-4)$. Дробный дрейф нуля в регистре кодирует ортогональный сдвиг осей в пространстве-времени через узлы \textbf{180} (экваториальная инверсия) и \textbf{10080} (минуты недели, связывающие инвариант октавы 28 и ядро 360: $28 \times 360 = 10080$). Глобальная адельная стабилизация Хассе гарантирует согласованность макро-календаря и виртуальных мотивов Гротендика во всех локальных точках.

\end{document}
